\section{Introduction}\label{sec:intro}

Retrieval-augmented generation (RAG) answers a question in two steps. A retriever selects passages from a corpus, and a language model, the reader, writes an answer from them \citep{lewis2020rag,izacard2021fid,ram2023incontext}. Teams that build such systems must choose the retriever. For dense retrieval this means choosing a text embedding model, and the usual guide is a retrieval benchmark such as BEIR \citep{thakur2021beir} or MTEB \citep{muennighoff2023mteb}. These benchmarks score ranked lists against relevance labels. They do not show how much of a retrieval gain survives the reader.

The question matters most for teams with little compute. A larger encoder costs more to run over a corpus and adds latency to every query. If a retrieval gain turns into a much smaller answer gain, a cheaper encoder may serve as well. Several studies report that retrieval improvements do not reliably improve answers \citep{salemi2024erag,mathur2026policy,medrano2026fusion,ghosh2026evidence}. Others find that retrieval metrics correlate well with generation quality when the two are aligned \citep{samuel2026beyond}. Readers are also sensitive to what else is in the context: passages that look relevant but lack the answer can lower accuracy \citep{cuconasu2024power,amiraz2025distracting,shi2023distracted}.

Earlier studies relate retrieval quality to answer quality through correlations across systems or topics \citep{salemi2024erag,samuel2026beyond,trappolini2025udcg}, through a ratio of answer F1 to recall across RAG methods \citep{peng2026neocorrag}, or through manipulations of a single retriever \citep{leto2024toward}. They do not report, for a set of interchangeable open encoders, the paired ratio of answer gain to retrieval gain with its uncertainty, per reader and dataset, next to the measured cost of each encoder. This study fills that gap for one well-defined setting.

We define a transfer ratio $T(R)$: the gain in exact match (EM) of retriever $R$ over a small baseline encoder, divided by its gain in Hit@5, both measured on the same queries. A ratio of 1 means every extra question with a relevant passage in the top five also gains a correct answer. A ratio near 0 means the retrieval gain does not reach the answer. We estimate $T(R)$ for BM25 and five dense encoders against MiniLM-L6, with two quantised 7--8B readers, on a single-hop dataset (Natural Questions) and a multi-hop dataset (HotpotQA). All runs used free Kaggle T4 GPUs.

We ask four research questions.
\begin{description}
\item[RQ1] How much of a Hit@5 gain over a small encoder reaches EM, and does this depend on the reader and the dataset?
\item[RQ2] Do larger or newer encoders give higher EM than a mid-size encoder (e5-base-v2)?
\item[RQ3] How do context choices change the answer: the number of passages, the kind of non-gold passages, the position of the gold passage, and reranking?
\item[RQ4] What does each encoder cost on free hardware, and do our sub-corpus and exact search distort the comparison?
\end{description}

The study ran in two phases. Phase~1 covered one reader on NQ. Before Phase~2, we registered six predictions (P1--P6) in a dated design document and then ran the new cells and experiments. We report each prediction as it came out, including the three that failed in at least one condition.

The contributions are as follows.
\begin{itemize}
\item A paired estimate of retrieval-to-answer transfer for six retrievers, two readers and two datasets, with bootstrap intervals and a reporting rule fixed in advance (Section~\ref{sec:res-transfer}).
\item An exact decomposition of the ratio into the hit--miss gap of the baseline and the change in EM given a hit or a miss. It separates the NQ pattern (low transfer) from the HotpotQA pattern (larger transfer driven by the second supporting passage) (Section~\ref{sec:res-decomp}).
\item Controlled tests of context composition and position, a reranker, top-$k$, approximate search and a full-corpus check, all run in the same pipeline (Sections~\ref{sec:res-context}--\ref{sec:res-cost}).
\item Measured encoding cost of each encoder on a free GPU, with five repetitions, and a full record of every deviation and failed run (Sections~\ref{sec:res-cost} and~\ref{sec:deviations}).
\end{itemize}

Section~\ref{sec:related} reviews related work. Sections~\ref{sec:method} and~\ref{sec:setup} describe the method and setup. Section~\ref{sec:results} reports the results, Section~\ref{sec:discussion} discusses them, and Section~\ref{sec:limits} lists the limits of the study.
